
In \Cref{thm_hs} and \Cref{sec_std_graphs}, we saw that the number of handshakes in a group of $n$ people or the number of edges of the complete graph $K_n$, is 
$$1 + 2 + 3 + \ldots + (n-1)  = \binom{n}{2} = \frac{n(n-1)}{2}.$$  
We often use sums in counting because cases add (\Cref{thm_cases_add}).

Similarly, in \Cref{sec_counting_steps_cases} and \Cref{sec_perms}, we saw that the number of ways to order $n$ distinct objects, or the number of permutations of the set $\{1, 2, 3, \ldots, n\}$, is
$$n(n-1)(n-2) \cdots(3)(2)(1)= n!$$
We often use products in counting because steps multiply (\Cref{thm_steps_multiply}).

In this chapter, we explore sums and products. We begin with two activities involving sums in \Cref{sec_model_sum_prod} {\secsumprod}, one of which introduces a representation of integers using sums of factorials.  We also introduce summation and product notation. \Cref{sec_binary_hex} {\secbinhex} continues our discussion of representing integers using sums, this time of binary, hexadecimal, and other base representations. Sums and products of an arbitrary number of integers are defined recursively, and so it is reasonable to investigate if there are explicit formulas for sums and products in \Cref{sec_sum_conj} {\secsumconj}.  We return to counting and polynomial algebra in \Cref{sec_arith_triangle} {\secarithtriangle} where we explore recursive and explicit formulas for $\binom{n}{k}$ as well as patterns in the Arithmetic triangle, some of which involve sums.  Finally, in  \Cref{sec_sum_proofs} {\secsumproofs} we look at techniques for proving a summation or product formula including proof by mathematical induction, combinatorial proof, and other proofs.